%% =========================================================================
%% Section 9: Conclusions and References
%% =========================================================================

\section{Conclusions and References}

\subsection{Summary of Empirical Discoveries}
This investigation conducted a survey-weighted life-course epidemiological evaluation of early-life disadvantage on later-life cognitive vitality and depressive symptoms among $66{,}606$ older adults in India using LASI Wave 1 microdata. Five core empirical conclusions emerge from this study:

\begin{enumerate}
    \item \textbf{The Enduring Affective Footprint of Childhood Adversity:}
    Childhood health shocks and economic poverty leave persistent psychological marks that adult socioeconomic mobility does not erase. Adults reporting poor childhood health had a \textbf{20\% higher prevalence} of depressive symptoms ($\text{APR} = 1.20$, 95\% CI: 1.07--1.34, $p = 0.002$), poor childhood SES conferred a \textbf{21\% higher prevalence} ($\text{APR} = 1.21$, 95\% CI: 1.03--1.42, $p = 0.019$), and childhood bedridden illness conferred a \textbf{16\% higher prevalence} ($\text{APR} = 1.16$, 95\% CI: 1.05--1.28, $p = 0.004$) after full adjustment for adult schooling and wealth. The Cumulative Adversity Index showed an \textbf{11.4\% risk increase per additional childhood hardship} ($p < 0.001$).
    
    \item \textbf{Cognitive Attenuation and the Equalizing Role of Education:}
    On an education-fair memory and orientation battery, the childhood poverty penalty ($\text{PR} = 1.61$ crude, $\text{APR} = 1.36$ demographic) \textbf{attenuated completely to the null} upon adjusting for adult education and wealth ($\text{APR} = 0.93$, 95\% CI: 0.82--1.05, $p = 0.241$). Formal schooling serves as a powerful protective resource, building cognitive reserve that buffers against clinical impairment regardless of childhood poverty.

    \item \textbf{The Social Mobility Asymmetry (Intellectual Resilience vs. Emotional Scars):}
    Our intergenerational mobility trajectory models revealed a profound life-course asymmetry: upwardly mobile adults who climbed from poor childhoods to top adult consumption quintiles experienced substantial cognitive recovery (impairment prevalence dropped from 13.9\% to 13.0\%, narrowing 60\% of the gap with affluent peers), but showed \textbf{no recovery in depressive symptoms} (prevalence remained elevated at 29.2\%, virtually identical to the 29.4\% among persistently poor peers). Adult affluence builds cognitive reserve but cannot heal the biological and psychological wounds of childhood precarity.

    \item \textbf{Gendered Vulnerability to Early Hardship:}
    Sex-stratified analyses revealed that older women disproportionately carry the psychological burden of childhood adversity ($\text{APR} = 1.23, p = 0.037$ for poor childhood SES among women vs. null $\text{APR} = 1.14, p = 0.146$ among men). This reflects the cumulative lifetime embodiment of gendered discrimination in nutrition, schooling, and domestic labor in mid-20th-century India.

    \item \textbf{Absence of Age Divergence and Oldest-Old Survivor Selection:}
    Joint Wald interaction tests confirmed that early disadvantage does not widen across age groups (all joint $p > 0.70$ for depression), refuting cumulative divergence. A nominal interaction for cognition ($p = 0.045$) was driven entirely by a mortality selection crossover at age 75+, where only exceptionally robust poor elders survived into advanced old age.
\end{enumerate}

\subsection{References and Bibliography}

\begin{enumerate}[leftmargin=1.0cm, label={[\arabic*]}]
    \item Arokiasamy, P., Bloom, D., Lee, J., Feeney, K., \& Ozolins, M. (2012). Longitudinal Aging Study in India: Vision, design, implementation, and preliminary findings. In \textit{Aging in Asia: Findings from New and Emerging Data Initiatives}. National Academies Press (US).
    \item Barker, D. J. P. (1990). The fetal and infantile origins of adult disease. \textit{BMJ: British Medical Journal}, 301(6761), 1111.
    \item Ben-Shlomo, Y., \& Kuh, D. (2002). A life course approach to chronic disease epidemiology: Conceptual models, empirical challenges and future directions. \textit{International Journal of Epidemiology}, 31(2), 285--293.
    \item Blau, P. M., \& Duncan, O. D. (1967). \textit{The American Occupational Structure}. New York: John Wiley \& Sons.
    \item Crimmins, E. M., Kim, J. K., Alley, D. E., Karlamangla, A., \& Seeman, T. (2005). Hispanic paradox in biological risk profiles: Latino individuals in the United States. \textit{American Journal of Public Health}, 95(7), 1305--1310.
    \item DiPrete, T. A., \& Eirich, G. M. (2006). Cumulative advantage as a mechanism for inequality: A review of theoretical and empirical developments. \textit{Annual Review of Sociology}, 32, 271--297.
    \item Elder, G. H. (1974). \textit{Children of the Great Depression: Social Change in Life Experience}. Chicago: University of Chicago Press.
    \item International Institute for Population Sciences (IIPS), National Programme for Health Care of the Elderly (NPHCE), MoHFW, Harvard T.H. Chan School of Public Health, \& University of Southern California (2020). \textit{Longitudinal Ageing Study in India (LASI) Wave 1, 2017--18, India Report}. Mumbai: IIPS.
    \item Kessler, R. C. (2003). Epidemiology of women and depression. \textit{Journal of Affective Disorders}, 74(1), 5--13.
    \item Khanal, S., \& Selvamani, Y. (2023). Early life disadvantage and depressive symptoms among older adults in India: Evidence from the Longitudinal Ageing Study in India. \textit{Journal of Affective Disorders}, 324, 452--460.
    \item Krieger, N. (2005). Embodiment: A conceptual glossary for epidemiology. \textit{Journal of Epidemiology \& Community Health}, 59(5), 350--355.
    \item Kuh, D., Ben-Shlomo, Y., Lynch, J., Hallqvist, J., \& Power, C. (2003). Life course epidemiology. \textit{Journal of Epidemiology \& Community Health}, 57(10), 778--783.
    \item Lee, J., Meijer, E., Phillips, D., \& Hu, P. (2021). Cognitive impairment in older adults in India: Evidence from the Harmonized Diagnostic Assessment of Dementia (DAD) study. \textit{Alzheimer's \& Dementia}, 17(S10), e053124.
    \item Lynch, S. M., \& Smith, H. L. (2005). Ageing, health, and selective survival: Conceptual and methodological issues. In \textit{Handbook of the Life Course} (pp. 373--392). Boston, MA: Springer.
    \item NITI Aayog (2023). \textit{Senior Care Reforms in India: Thinking of the Future}. Government of India, New Delhi.
    \item Patel, V., Kirkwood, B. R., Pednekar, S., Weiss, H., \& Mabey, D. (2006). Risk factors for common mental disorders in women: Population-based longitudinal study in Goa, India. \textit{The British Journal of Psychiatry}, 189(6), 547--555.
    \item Sorokin, P. A. (1927). \textit{Social Mobility}. New York: Harper \& Brothers.
    \item Stern, Y. (2002). What is cognitive reserve? Theory and research application of the reserve concept. \textit{Journal of the International Neuropsychological Society}, 8(3), 448--460.
    \item Vaupel, J. W., \& Yashin, A. I. (1985). Heterogeneity's ruses: Some surprising effects of selection on population dynamics. \textit{The American Statistician}, 39(3), 176--185.
    \item Zou, G. (2004). A modified Poisson regression approach to prospective studies with binary data. \textit{American Journal of Epidemiology}, 159(7), 702--706.
\end{enumerate}
